Con los resultados obtenidos se puede concluir que el comportamiento práctico de los algoritmos está fuertemente relacionado con su complejidad teórica y su implementación en el caso de los algoritmos de ordenamiento, ya que Patience Sort alcanzó aproximadamente 16 segundos para $n=10^5$ en el caso aleatorio D1, mostrando un crecimiento cercano a $O(n^2)$ en situaciones desfavorables. Luego se pudo notar que Merge Sort, Quick Sort y Sort presentaron tiempos bastante menores, cercanos a un comportamiento $O(n\log n)$. Además, en los algoritmos que usan más estructuras auxiliares, se pudo ver un mayor consumo de memoria, como Patience Sort, que alcanzó cerca de 2000 KB. En los algoritmos de multiplicación de matrices, Strassen mostró alrededor de 370 segundos para $n=2^{10}$, que al momento de compararlo con Naive, se ve una gran diferencia ya que Naive demoró aproximadamente 8 segundos, demostrando que una mejor complejidad asintótica no siempre conlleva a un menor tiempo de ejecución para los tamaños analizados en el caso de multiplicación de matrices. Finalmente, se cumplió el objetivo de relacionar el comportamiento experimental de los algoritmos con su complejidad teórica.
